% Folder 84, batch 01 (archivists' pages 1-20).
% Pass: Opus 5.5 (claude-opus-5-5), 2026-10-09 — first pass, unchecked against the pages by a human.
% Revised: Opus 5.5 (claude-opus-5-5), 2026-10-10 — re-read on the 700 dpi tiles; 35 \ill and 18 \uncertain resolved, 59 remain.
%
% Pages 4, 6 and 7 are unrelated versos (typescript pages, a computer
% listing) and are absent. Page 2 carries, below the calculations, a short
% personal note to someone, not transcribed.
\documentclass[11pt,a4paper]{article}
\input{../preamble/grothendieck}

\folder{84}
\batch{1}
\pages{1}{20}
\foldertitle{[Formes quadratiques] : notes manuscrites (s.d.).}
\dating{[à partir de 1982-vers 1986]}
\watermark{Édition de démonstration}

\begin{document}


\page{1}
\note{feuillet de calculs très chargé, en plusieurs couches ; on transcrit les formules dans l'ordre où elles se lisent, en comprimant la disposition}

$R_1 \rightleftarrows R_0$, avec $\beta : R_1 \to R_0$ et $\alpha : R_0 \to R_1$. \struck{hom. de monoïdes}

$R_0$, $R_1$ monoïdes ;

$\alpha : R_0 \to R_1$ hom. de monoïdes, $(L, \delta \in \Gamma(L^{\otimes 2}), T_0 \in \Gamma L_0) \longmapsto (L, \Gamma)$\note{la dernière lettre se lit $\Gamma$ ; on attendrait $\delta$}

$\beta : R_1 \to R_0$ \uncertain{multiplicat.}, \ill{} (\uncertain{non compat. avec} produits), $(L,\delta) \longmapsto (L, 4\delta, 0)$.

$R_* = R_0 \sqcup R_1$, loi de composition $(x,y) \mapsto x \wedge y$ (comm., ass., unitaire), compat. avec $R_* \to \mathbb{Z}/2\mathbb{Z}$, i.e.
\[
  R_0 . R_0 \subset R_0, \qquad R_1 R_1 \subset R_0, \qquad R_1 . R_0 = R_0 . R_1 \subset R_1 .
\]
\begin{align*}
  x, y \in R_0 &\Longrightarrow x \wedge y = x . y \\
  x, y \in R_1 &\Longrightarrow x \wedge y = \beta(x.y) \qquad (\ldots : \beta(x)\beta(y)) \\
  x \in R_1,\ y \in R_0 &\qquad x \wedge y = y \wedge x = x\, \alpha(y)
\end{align*}
\note{dans la dernière ligne, une expression biffée et noircie précède $x\,\alpha(y)$ ; à la deuxième, la parenthèse se lit « (\uncertain{pas néc. égal} : $\beta(x)\beta(y)$) »}

$(x \wedge y) \wedge z = x \wedge (y \wedge z)$ :
\begin{itemize}
  \item a) $x, y, z \in R_0$ : OK (\struck{\ill{}}) — ass. dans $R_0$ ;
  \item b) $x \in R_1$, $y, z \in R_0$ : $(x . \alpha(y))\alpha(z) = x . \alpha(y.z)$ — multipl. de $\alpha$ + ass. dans $R_1$ ;
  \item c) $y \in R_1$, $x, z \in R_0$ : $(\alpha(x) y)\alpha(z) = \alpha(x)(y\,\alpha(z))$ — ass. dans $R_1$ ;
  \item d) $z \in R_1$, $x, y \in R_0$ : $\alpha(x.y).z = \alpha(x)(\alpha(y).z)$ — multipl. de $\alpha$ + ass. de $R_1$ ;
  \item b') $x \in R_0$, $y, z \in R_1$ : $\beta((\alpha(x) y).z) \overset{?}{=} x . \beta(y.z)$ \quad $(*)$ ;
  \item c') $y \in R_0$, $x, z \in R_1$ : $\beta((x\,\alpha(y)).z) = \beta(x.(\alpha(y).z))$ — ass. dans $R_1$ ;
  \item d') $z \in R_0$, $x, y \in R_1$ : $\beta(x.y)\, z \overset{?}{=} \beta(x(y\,\alpha(z)))$ — comme $(*)$ + comm. et ass. dans $R_1$ ;
  \item a') $x, y, z \in R_1$ : $\alpha(\beta(x.y)).z \overset{?}{=} x\, \alpha(\beta(y.z))$.
\end{itemize}
\note{sous d') il récrit les deux membres : $z\,\beta(y.x)$ et $\beta((\alpha(z)y)x)$}

Donc il faut vérifier

① $\beta(\alpha(x).yz) = x\,\beta(yz)$, i.e. $\beta(\alpha(x)u) = x\,\beta(u)$ ; $\beta(\alpha(x)) = x.\beta(1)$ (cas partic. $y = z = 1$) ;

② $(\alpha\beta(xy))z = x\,\alpha\beta(yz)$ ; $(\alpha\beta(x))z = x(\alpha\beta(z))$ si $y = 1$ ; et si $z = 1$ : $\alpha\beta(x) = x.\alpha\beta(1) = x\xi$.

Et ② devient donc $(xy\xi)z = x(yz\xi)$, OK par ass. et comm. dans $R_1$.

\[
  \text{(A)}\quad \beta(\alpha(x)u) = x\,\beta(u), \qquad
  \text{(B)}\quad \alpha\beta(x) = x.\eta, \quad \eta = \alpha\beta(1) = \alpha(\zeta), \quad \zeta = \beta(1)
\]
avec $\zeta = (\underline{O}_X, 4\,1_{\underline{O}}, 0)$, d'où
\[
  \beta\alpha(x) = x\zeta \quad \text{OK}, \qquad (\beta x)(\beta y) = \zeta\, \beta(x.y).
\]
\struck{$\beta(1)\beta(x.y) = \ldots$} \quad $\beta(x)\beta(y) \overset{?}{=} \zeta . \beta(x.y)$, via (A) et (B) : $\beta(\alpha\beta(x)\,y) = \beta(\alpha(\zeta)\,x\,y) = \zeta\,\beta(xy)$.
\note{la chaîne d'égalités sous la boîte est reconstituée de plusieurs fragments reliés par des flèches ; l'ordre est conjectural}

Vér. les \ill{} d'\uncertain{espèces} : $x = (L, \delta)$, $\beta(x) = (L, 4\delta, 0)$, $\alpha\beta(x) = (L, 4\delta) = (L,\delta)(\underline{O}, 4\,1_X)$ ;
$x = (L, \delta, T_0)$, $u = (L', \delta')$ ; $\alpha(x) = (L, \delta)$, $\alpha(x).u = (LL', \delta\delta')$, $\beta(\alpha(x).u) = (LL', 4\delta\delta', 0)$ ;
$\beta(u) = (L', 4\delta', 0)$, $x\,\beta(u) = (LL', 4\delta\delta', 0)$.

\marginal{encadré : $R_* = R \times_{\mathbb{Z}} \mathbb{Z}/2\mathbb{Z}$. On définit $w_n : R_* \to R_0$ par : a) si $n = 2\nu$, $w_n(\xi) = x\,\zeta^{\nu}$ ; b) si $n = 2\nu + 1$, $w_n : R_1 \to R_0$, $w_n(x) = \beta(x)\,\zeta^{\nu}$. On trouve ainsi un hom. $w : R_* \to R_0$. Un \ill{}\note{le bord de la feuille est déchiré} hom. $w' : R_* \to \mathbb{Z} \times R_0$.}

$E = \bigoplus_{1 \le i \le n} (L_i, \delta_i)$,
\[
  w_1(E) = \bigwedge_{1 \le i \le n} (L_i, 4\delta_i, 0) = \Bigl(\bigotimes L_i,\ 4^{n} \textstyle\prod \delta_i,\ 0\Bigr),
\]
\[
  \mathrm{Arf}(E) = \Bigl(\bigotimes L_i,\ 2^{\nu} \textstyle\prod \delta_i,\ 0\Bigr),
  \qquad \mathrm{Arf}(E) = w_1(E).\,\xi^{\ldots}
\]
avec $\nu = n$ si $n$ pair, $\nu = n+1$ si $n$ impair ; $\nu = 2\nu'$, $\nu' = \bigl[\tfrac{n+1}{2}\bigr]$ ; $\xi = (\underline{O}_X, 2^{?}1, 1)$.

$w_1(L_n) = (L, 4\delta, 0)$, \quad $\mathrm{Arf}(L) = (L, 4\delta, 0)$.

\struck{On peut prendre} NB si $2$ inv. dans $X$ : \struck{$\ldots \simeq \ldots$} \struck{$\mathrm{Arf}(E) \simeq \ldots$}

On a
\[
  c_1(E) = \mathrm{Arf}(E).\,\xi^{(n - \nu') = \nu''}, \qquad \nu'' = \begin{cases} \tfrac{n}{2} & \text{si } n \text{ pair} \\ \tfrac{n-1}{2} & \text{si } n \text{ impair} \end{cases}
\]
où $\xi = (\underline{O}_X, 4, 0)$, $\eta = \mathrm{Arf}(\mu_{2X}/X)$.
\note{à côté, entre deux traits : $c_1(E) \simeq (\bigotimes L_i, \delta, 0)$, $\mathrm{Arf}(E) \simeq (\bigotimes L_i, \delta, 0)$ ; et une remarque : « donc $\mathrm{Arf}(E)$ \uncertain{est un invariant} plus fin que $c_1(E)$ (si $2$ \uncertain{inversible} sur $X$) »}

Mais $c_1(E)$ \ill{}
\note{la page s'arrête sur ces mots}

\page{2}
$R_1$, $R_0$ ;
\[
  R_1 \times R_1 \to R_0, \quad R_0 \times R_0 \to R_0, \quad R_0 \times R_1 \to R_1, \quad R_1 \times R_0 \to R_1 .
\]
$R = R_1 \sqcup R_0$ est un groupoïde unitaire ; \struck{\ill{}} $R \xrightarrow{\varepsilon} \mathbb{Z}/2\mathbb{Z}$ est un hom. de monoïdes ;
$R_1 = \varepsilon^{-1}(1)$, $R_0 = \varepsilon^{-1}(0)$ ; $R_{00} \subset R_0$, partie \ill{} $= H^1(X, \mathbb{Z}/2\mathbb{Z})$.

$R_1 \to R_0$ n'est \emph{pas} un hom. de monoïde ? \quad $(L,\delta) \mapsto (L, 4\delta, 0)$\note{sous $(L,\delta)$, une flèche vers le bas et deux mots : \ill{}}.

\emph{Groupoïde} $(L, \delta, T_0)$ ; $(L,\Gamma) \mapsto (\pi\ \omega)(L, \delta, T_0) \mapsto (L,\delta)$ ; $(L, 4\delta, 0)$.

\emph{Paires} $(L, \delta)$ ; \quad $(L, \delta)$, $L$.

$R_1 \to H^1(L)$\note{lecture douteuse}

$\theta = 2\delta$, \quad $\theta' = 2\delta'$, \quad $\theta\theta' = 4\delta\delta'$, \quad $\theta\delta' = 2$ ; \quad $\theta = 2\delta$.

\begin{align*}
  (L, \delta)(L', \delta') &= (LL', 4\delta\delta', 0) \\
  (L, \delta)(L', \delta', T'_0) &= (LL', 2\delta\delta') \\
  (L, \delta, T_0)(L', \delta') &= (LL', \delta\delta') \\
  (L, \delta, T_0)(L', \delta', T'_0) &= (LL', \delta\delta', T_0 T'_0)
\end{align*}
\note{dans la deuxième ligne le $2$ est surchargé et douteux}

« deux ou trois produits » \ill{} : $(LL'L'', 4\delta\delta'\delta'', 0 \ldots)$, $(LL'L'', \delta\delta'\delta'', T_0T'_0T''_0)$ \ill{}.

(\emph{Commutation})


\page{3}
$A_1$, $A_2$, $A_3$, de bases $1, u_1$ ; $1, u_2$ ; $1, u_3$, avec $u_i^2 + e_i = 0$ ;
\[
  A_1 \otimes A_2 \otimes A_3 \overset{\mathrm{df}}{=} A .
\]
Base : $1$, $u_1$, $u_2$, $u_3$, $u_2u_3$, $u_3u_1$, $u_1u_2$, $u_1u_2u_3$.

extensions quadratiques, involutions canoniques $\sigma_1, \sigma_2, \sigma_3$ ($\sigma_i(u_i) = -u_i$) se prolongent à $A$ de façon canonique par $\sigma_1 \otimes \mathrm{id} \otimes \mathrm{id}$ etc.

\[
\begin{array}{lll}
  k & 1 & \sigma_1, \sigma_2, \sigma_3, \sigma_2\sigma_3, \sigma_3\sigma_1, \sigma_1\sigma_2, \sigma_1\sigma_2\sigma_3 = \tau \\
  A_1 & 1, u_1 & \sigma_2, \sigma_3, \sigma_2\sigma_3 \\
  A_2 & 1, u_2 & \sigma_3, \sigma_1, \sigma_3\sigma_1 \\
  A_3 & 1, u_3 & \sigma_1, \sigma_2, \sigma_1\sigma_2 \\
  B_1 & 1, u_2u_3 & \sigma_1, \sigma_2\sigma_3, \sigma_1\sigma_2\sigma_3 = \tau \\
  B_2 & 1, u_3u_1 & \sigma_2, \sigma_3\sigma_1, \sigma_2\sigma_3\sigma_1 = \tau \\
  B_3 & 1, u_1u_2 & \sigma_3, \sigma_1\sigma_2, \sigma_3\sigma_1\sigma_2 = \tau \\
  B_0 & 1, u_1u_2u_3 & \sigma_2\sigma_3, \sigma_3\sigma_1, \sigma_1\sigma_2
\end{array}
\]
\[
  B_1 \simeq A_2 \wedge A_3, \quad B_2 \simeq A_3 \wedge A_1, \quad B_3 \simeq A_1 \wedge A_2, \quad B_0 \simeq A_1 \wedge A_2 \wedge A_3 .
\]
\note{la colonne de droite du tableau donne les involutions qui laissent fixe la sous-algèbre de la ligne}

\[
\begin{array}{lll}
  A'_1 = (A_2, A_3, B_1) & 1, u_2, u_3, u_2u_3 & \sigma_1 \\
  A'_2 = (A_3, A_1, B_2) & 1, u_3, u_1, u_3u_1 & \sigma_2 \\
  A'_3 = (A_1, A_2, B_3) & 1, u_1, u_2, u_1u_2 & \sigma_3 \\
  B'_1 = (A_1, B_1, B_0) & 1, u_1, u_2u_3, u_1u_2u_3 = u & \sigma_2\sigma_3 \\
  B'_2 = (A_2, B_2, B_0) & 1, u_2, u_3u_1, u_2u_3u_1 = u & \sigma_3\sigma_1 \\
  B'_3 = (A_3, B_3, B_0) & 1, u_3, u_1u_2, u_3u_1u_2 = u & \sigma_1\sigma_2 \\
  B'_0 = (B_1, B_2, B_3) & 1, u_2u_3, u_3u_1, u_1u_2 & \sigma_1\sigma_2\sigma_3
\end{array}
\]
correspondant aux sept droites vect. de $\mathbb{F}_2^{\,3}$ \struck{\ill{}}, i.e. aux \struck{sept} sept éléments non nuls.

\marginal{correspondant aux sept hyperplans de $\mathbb{F}_2^{\,3}$, i.e. \ill{} \ill{} — \struck{aux sept droites de $\mathbb{F}_2^{\,3}$ i.e. aux} sept éléments non nuls}

$A_i \otimes A'_i \xrightarrow{\ \sim\ } A$ \quad [$B_i \subset B'_i$] \quad $B_0 \otimes B'_0 \xrightarrow{\ \sim\ } A$, si les \uncertain{$e_i$ inv.} i.e. les $A_i$ étales.

\begin{align*}
  A_1 &\subset A'_2, A'_3, B'_1 &  B_1 &\subset A'_1, B'_1, B'_0 \\
  A_2 &\subset A'_3, A'_1, B'_2 &  B_2 &\subset A'_2, B'_2, B'_0 \\
  A_3 &\subset A'_1, A'_2, B'_3 &  B_3 &\subset A'_3, B'_3, B'_0 \\
      &                          &  B_0 &\subset B'_1, B'_2, B'_3
\end{align*}
\note{sous la dernière ligne : $B_0(A_1)$, $B_0(A_2)$, $B_0(A_3)$}

il y a $7 + 7 = 14$ situations « biquadratiques » typiques (quand les $A_i$ sont étales…)


\page{5}
$1, u$ \quad $\longmapsto -u + \lambda$

$1, u \longmapsto \bar{u} + \alpha$
\[
  N(\bar u + \alpha) = N(\bar u) \quad (= N(u)), \qquad
  \mathrm{Tr}(\bar u + \alpha) = \mathrm{Tr}\,\bar u \quad (= \mathrm{Tr}\,u)
\]
\[
  N(\bar u) + \alpha\,\mathrm{Tr}(\bar u) + \alpha^2 = N(\bar u)
\]
\[
  \begin{cases} \alpha(\mathrm{Tr}\,\bar u + \alpha) = 0 & (\mathrm{Tr}\,\bar u = b) \\ 2\alpha = 0 \end{cases}
  \Longrightarrow \alpha = 0 \text{ si } 2 \text{ inv.}
\]

\struck{$\lambda + \mu u$} \qquad $u\bar u = N(u)$
\[
  N(v) = 1 \Longrightarrow v = u\,\bar u^{-1} \;(= u^2 N(u)^{-1}), \qquad u^2 + bu + c = 0, \quad \mathrm{Tr}\,u = -b .
\]
\begin{align*}
  v &= \alpha + \beta u, & N(v) &= \alpha^2 - \alpha\beta b + \beta^2 c = 1 \\
  \bar v &= (\alpha - \beta b) - \beta u, & &= \alpha(\alpha - \beta b) \quad \text{si } c = 0 \\
  v \bar v &= \alpha(\alpha - \beta b) - \beta^2 c \ldots
\end{align*}
\note{dans $N(v)$ le signe devant $\alpha\beta b$ est noirci}
\[
  v = \alpha + \beta u, \quad \bar v = \alpha + \beta \bar u, \qquad
  v\bar v = \alpha^2 + \beta^2 + \alpha\beta\,(u + \bar u)
\]
\note{sous $(u + \bar u)$ : $\mathrm{Tr}\,u$ ; ce second calcul est biffé en partie}

$w = \lambda 1 + \mu u$ \& \struck{$N(u) = 0$} $N(u) = 0$, donc \struck{$N(w)$ inv.} $N(w)$ inv. $\Longrightarrow$ $\lambda$ inv.
\[
  N(w) = \lambda^2 - \lambda\mu b = \lambda(\lambda - \mu b) ; \qquad \alpha(\alpha - \beta b) = 1 .
\]

\begin{tikzcd}
  & SW/\mu_2 = W^*/G_m \arrow[r] & 1 & \\
  1 \arrow[r] & SW \arrow[r] \arrow[u] & W^* \arrow[r, "N"] & G_m \arrow[r] & 1 \\
  1 \arrow[r] & \mu_2 \arrow[r] \arrow[u] & G_m \arrow[r, "{(\ )^2}"] \arrow[u] & G_m \arrow[r] \arrow[u, no head] & 1
\end{tikzcd}

\note{la flèche $G_m \to W^*$ porte « cas \ill{} » ; la flèche $G_m \to G_m$ de la dernière ligne porte un exposant $2$ de lecture douteuse ; les deux $G_m$ de droite sont reliés par un double trait ; la première ligne est en partie biffée et des flèches verticales partent encore vers le bas de la page}

\section{Formes quadratiques binaires et revêtements quadratiques}
\note{titre souligné, en tête de la p.~8}

\page{8}
1. Forme quadratique à scalaire près, structure de similitude.

1.1. $(X, \underline{O}_X)$ désigne un schéma, plus généralement un topos loc. annelé. Soit $E$ un Module sur $X$,
\[
  \underline{\mathrm{Quad}}(E) \simeq \underline{\mathrm{Hom}}(\Gamma^2(E), E)
\]
\note{la seconde lettre se lit $E$, où l'on attendrait $\underline{O}_X$}
le faisceau des formes quadratiques sur $E$. Si $E$ est loc. libre de type fini, on a un iso. can.
\[
  \underline{\mathrm{Quad}}(E) \simeq \underline{\mathrm{Sym}}^2(E^\vee),
\]
(en général on a seulement un homomorphisme can.
\[
  \underline{\mathrm{Sym}}^2(E^\vee) \longrightarrow \underline{\mathrm{Quad}}(E),
\]
qui transforme $f.g$ ($f, g \in \Gamma E^\vee$) dans $Q_{fg} : x \mapsto f(x)g(x)$ sur $E$…).
Le faisceau \struck{\ill{}} en groupes $\underline{O}_X^*$ opère sur $\underline{\mathrm{Quad}}(E)$, d'où un faisceau quotient
\[
  \underline{\mathrm{Sim}}(E) \simeq \underline{\mathrm{Quad}}(E)/\underline{O}^* .
\]
Une section du faisceau quotient s'appelle une \emph{structure de similitude} sur $E$. Une telle section correspond à un sous-faisceau ensembliste de $\underline{\mathrm{Quad}}(E)$, qui \ill{} un sous-faisceau \struck{loc} soit : fibres non vides et sur lequel $\underline{O}^*$ opère transitivement\note{les lignes, inclinées, se chevauchent ici ; l'ordre des membres de phrase est incertain}. Localement, une telle structure est définie par une section $Q$ de $\underline{\mathrm{Quad}}(E)$, i.e. par une forme quadratique sur $E$, deux telles sections $Q$, $Q'$ définissant la même structure de similitude ssi $\exists$ localement $a \in \Gamma(\underline{O}_X^*)$ avec $Q' = aQ$. Le cas le plus intéressant est celui où \struck{$a \mapsto aQ$}
\[
  a \longmapsto a.Q : \underline{O}_X \ (\text{sans } * \,!) \longrightarrow \underline{\mathrm{Quad}}(E)
\]
est injectif. Dans ce cas, si $Q'$ \struck{définit la même} définit la même

\page{9}
structure de similitude ssi $\exists\, a \in \Gamma \underline{O}_X^*$ (globalement) telle que $Q' = aQ$, et cette $a$ est alors unique (ce qui assure l'unicité). \struck{On appelle des} \add{structure de similitude} \emph{fidèle} : définie loc. par une $Q$ telle que $a \mapsto aQ$ injectif \struck{(ce qui est indépendant du choix des $Q$, seulement de la structure de similitude)}.

1.2. \struck{Soit} Pour $P \subset \underline{\mathrm{Quad}}(E)$ une structure de similitude, on peut regarder le sous-Module engendré, soit
\[
  \Delta = \Delta_P \subset \underline{\mathrm{Quad}}(E),
\]
et on voit que
\[
  P \longmapsto \Delta_P
\]
donne une bijection entre l'ens. des structures de similitudes sur $E$, et l'ens. des sous-Modules de $\underline{\mathrm{Quad}}(E)$ qui sont localement monogènes. \struck{Les formes} Les structures \emph{fidèles} de similitude fidèles correspondent aux sous-modules
\[
  \Delta \hookrightarrow \underline{\mathrm{Quad}}(E)
\]
qui sont inversibles. Un \uncertain{hom.} quelc. $\Delta \to \underline{\mathrm{Quad}}(E)$ définit évidemment une structure \ill{}.

\emph{NB} Dans le cas où \ill{} $E$ est loc. \ill{} de type fini, dans \uncertain{ce cas} il \ill{}

1.3. \ill{} \ill{} \ill{} $\underline{\mathrm{Quad}}(E) \simeq \underline{\mathrm{Sym}}^2(E)$. Une structure \ill{} de similitudes est dite \emph{strictement fidèle} (ou universellement fidèle) si l'homom. précédent est « universellement injectif », ou encore (ce qui revient au même) injectif sur chaque fibre résiduelle.
\note{ce passage est très surchargé : une note interlinéaire précédée de « NB » et le numéro 1.3 s'y enchevêtrent ; la formule se lit $\underline{\mathrm{Sym}}^2(E)$ là où 1.1 donnait $\underline{\mathrm{Sym}}^2(E^\vee)$}

\page{10}
Les structures de similitude strictement fidèles correspondent donc aux \emph{sous-fibrés vectoriels de rang 1} du fibré vectoriel $\mathbb{W}(\underline{\mathrm{Quad}}(E))$], ou encore aux sections du fibré en espaces projectifs $\mathbb{P}^\vee(\underline{\mathrm{Quad}}(E))$.

1.4. Soient $\Delta$ un Module inv. sur $X$, et
\[
  \Delta \xrightarrow{\ i\ } \underline{\mathrm{Quad}}(E)
\]
un hom. (pas nécessairement injectif), définissant \add{donc} une structure de similitude $\Sigma$ sur $E$. Je dis qu'on a une forme quadratique \add{$Q_\Sigma$ ou $Q_i$} (canonique) à valeurs dans le module inversible $\Delta^\vee$. Pour la définir, je vais d'abord définir (indép. du choix de $i$ ou de $\Sigma$) une forme quadratique canonique
\[
  Q = Q_E : E \longrightarrow \underline{\mathrm{Quad}}(E)^\vee :
\]
\struck{elle associe} sa valeur, pour une section locale $x$ de $E$, est la forme linéaire
\[
  Q(x) = \{ f \longmapsto f(x) \} \qquad \text{i.e.} \quad f(x) = \langle f, Q(x) \rangle
\]
\struck{(pour toute section (locale) $f$ de $\underline{\mathrm{Quad}}(E)$, $x$ de $E$.} \struck{\ill{} $\underline{\mathrm{Quad}}(E)$).} Celle-ci dépend de $x$ de façon quadratique, i.e. la fonction $E \times E \to \underline{\mathrm{Quad}}(E)^\vee$,
\[
  (x, y) \longmapsto Q(x + y) - Q(x) - Q(y) = \{ f \longmapsto \underbrace{f(x+y) - f(x) - f(y)}_{\varphi_f(x,y)} \}
\]

\page{11}
est bilinéaire en $x, y$. Comme on a, par transposition de $i : \Delta \hookrightarrow \underline{\mathrm{Quad}}(E)$,
\[
  \underline{\mathrm{Quad}}(E)^\vee \xrightarrow{\ i^\vee\ } \Delta^\vee ,
\]
on en déduit une forme quadratique composée
\[
  Q_i = Q_\Sigma : E \longrightarrow \Delta^\vee .
\]
\note{au-dessus de la flèche, la composition est dessinée : $Q : E \to \underline{\mathrm{Quad}}(E)^\vee$ puis $i^\vee$}
Dans le cas d'une structure de similitude $\Sigma$ \add{fidèle}, $(\Delta, i)$ ne dépend, à iso. \add{can.} unique près, que de $\Sigma$, et il en est donc de même de $Q_i$, ce qui justifie la notation $Q_\Sigma$. \struck{Quand}

\emph{NB} Quand on se donne une base $e$ de $\Delta$, d'où $\Delta \simeq \Delta^\vee \simeq \underline{O}$, alors $Q_\Sigma$ s'identifie à une forme quadr. ordinaire $E$ à val. dans $\underline{O}$, et on a \ill{} \struck{\ill{}} \add{$f = i(e)$}.

\marginal{1.5. La donnée d'une structure de similitudes (\emph{fidèle} \add{strictement}) sur $E$ \uncertain{équivaut} à celle d'un Module inv. $\Delta$ sur $X$ et d'une forme quadratique $Q : E \to \Delta^\vee$ (\ill{} $Q(E)$ engendre $\Delta^\vee$, i.e. \struck{$Q$} \uncertain{sur aucune} fibre résiduelle identiquement nulle).}

2.1. On suppose à partir de maintenant que $E$ est loc. libre de rang 2. Soit
\[
  L = \underline{\det} E = \textstyle\bigwedge^2 E
\]
le Module déterminantiel, qui est un Module inversible. Notons l'iso. can.
\[
  E \xrightarrow{\ \sim\ } \check E \otimes L
\]
donné par $\bigl(x \longmapsto (y \longmapsto x \wedge y)\bigr)$, i.e. provenant de l'accouplement canonique
\[
  E \times E \longrightarrow L = \textstyle\bigwedge^2 E : \quad (x, y) \longmapsto x \wedge y .
\]

\page{12}
On en déduit des iso. can.
\[
  \begin{cases}
    \check E \simeq E \otimes L^\vee \\
    \underbrace{\underline{\mathrm{Sym}}^2(E^\vee)}_{\underline{\mathrm{Quad}}(E)} \simeq \underbrace{\underline{\mathrm{Sym}}^2(E)}_{\underline{\mathrm{Quad}}(E^\vee)} \otimes \underbrace{L^{\vee \otimes 2}}_{L^{\otimes(-2)}}
  \end{cases}
\]
i.e. $\underline{\mathrm{Quad}}(E) \simeq \underline{\mathrm{Quad}}(\check E) \otimes L^{\otimes -2}$.
Il en résulte un isom. canonique de faisceaux
\[
  \underline{\mathrm{Sim}}(E) \xrightarrow[\ \sim\ ]{\ \mathrm{can}_E\ } \underline{\mathrm{Sim}}(\check E)
\]
qui transforme le sous-Module loc. monogène $\Delta$ de $\underline{\mathrm{Sym}}^2(E^\vee) \simeq \underline{\mathrm{Quad}}(E)$ dans
\[
  \Delta \otimes L^{\otimes 2} \subset \underline{\mathrm{Quad}}(\check E) \otimes L^{\otimes 2}
  \simeq \underline{\mathrm{Quad}}(\check E) \otimes L^{\otimes(-2)} \otimes L^{\otimes 2}
  \simeq \underline{\mathrm{Quad}}(\check E).
\]
\note{avant $\underline{\mathrm{Quad}}(\check E) \otimes L^{\otimes 2}$, un $\underline{\mathrm{Sym}}^2(E)$ est biffé}
On trouve que $\mathrm{can}_E$ et $\mathrm{can}_{\check E}$ sont inverses l'un de l'autre. Les structures de similitude \emph{fidèles}, strictement fidèles, non dégénérées (\uncertain{universellement}…) pour $E$ et pour $E^\vee$ se correspondent.

2.2. Considérons \add{encore} l'iso. can. (dépendant du choix d'un signe)
\[
  \check E \simeq E \otimes L^\vee ,
\]
en \add{$\otimes$-}multipliant à droite par $E$, on trouve
\[
  \check E \otimes E \simeq E \otimes E \otimes L^\vee .
\]
D'autre part \struck{(\ill{} $\underline{O}_X$-Module)} on a une suite exacte canonique (qui ne fait intervenir aucun choix de signe…) :
\[
  0 \to \textstyle\bigwedge^2 E \longrightarrow E \otimes E \longrightarrow \underline{\mathrm{Sym}}^2(E) \to 0 ,
  \qquad x \wedge y \longmapsto x \otimes y - y \otimes x
\]

\page{13}
et dans le cas actuel où $\bigwedge^2 E \overset{\mathrm{df}}{=} L$ est inversible, \struck{et \ill{}} donc
\[
  0 \to L \longrightarrow E \otimes E \longrightarrow \underline{\mathrm{Sym}}^2(E) \to 0 ,
\]
d'où en multipliant par $L^\vee$
\[
  0 \to \underline{O} \longrightarrow (E \otimes E) \otimes L^\vee \longrightarrow \underline{\mathrm{Sym}}^2(E) \otimes L^\vee \to 0
\]
avec $(E \otimes E) \otimes L^\vee \simeq E^\vee \otimes E \simeq \underline{\mathrm{End}}(E)$.

Je dis que par l'iso. can. vertical, le sous-Module $L \otimes L^\vee \simeq \underline{O}$ est transformé en le sous-module $\underline{O}.1_E$ de $\underline{\mathrm{End}}(E)$, de façon plus précise que l'on a un diagramme \struck{commutatif des} d'isomorphismes

\begin{tikzcd}[column sep=small, nodes={font=\small}]
  0 \arrow[r] & L \otimes L^\vee \arrow[r] \arrow[d, no head, "\wr"] & (E \otimes E) \otimes L^\vee \arrow[r] \arrow[dd, "\wr"] & \underline{\mathrm{Sym}}^2(E) \otimes L^\vee \arrow[r] \arrow[dd, "\wr"] & 0 \\
  & \underline{O} \arrow[d, "{\mathrm{id}_{\underline{O}}}", "\wr"'] & & & \\
  0 \arrow[r] & \underline{O} \arrow[r] & \underline{\mathrm{End}}(E) \arrow[r] & \underline{\mathrm{End}}(E)/\underline{O} \arrow[r] & 0
\end{tikzcd}

\note{au-dessus de $L$ : $\det E$ ; au-dessus de $\underline{\mathrm{Sym}}^2(E) \otimes L^\vee$ : $\underline{\mathrm{Quad}}(\check E) \simeq \underline{\mathrm{Quad}}(E) \otimes L^{\otimes 2}$ ; sous $\underline{\mathrm{End}}(E)$ : $\simeq \check E \otimes E \simeq E \otimes L^\vee$ \uncertain{$\otimes E$}}

d'où par \ill{} l'\uncertain{isom.} \ill{}
\[
  \underline{\mathrm{Sym}}^2(E) \otimes L^\vee \simeq \underline{\mathrm{Quad}}(E) \otimes L \qquad (\simeq \underline{\mathrm{Quad}}(E, L))
\]
\[
  \wr\!\downarrow
\]
\[
  \underline{\mathrm{End}}(E)/\underline{O}
\]

\page{14}
On en conclut : \emph{les structures de similitude sur les Modules (loc. libres de rang 2) $E$ correspondent aux sous-Modules loc. monogènes} \add{Modules (loc. libres de rang 2)} \emph{des} $\underline{\mathrm{End}}(E)/\underline{O}$, i.e. aux sous-Modules $A$ de $\underline{\mathrm{End}}(E)$ tels que
\begin{itemize}
  \item 1°) $A \supset \underline{O}.1_E$
  \item 2°) $A/\underline{O}.1_E$ soit loc. monogène.
\end{itemize}
Les structures de similitude strictement fidèles correspondent aux sous-fibrés vectoriels \add{de rang 2} (de $\underline{\mathrm{End}}(E)$) qui contiennent le sous-fibré $\underline{O}.1_E$.

\emph{Lemme 1} Soit $A$ un sous-Module de $\underline{\mathrm{End}}(E)$ qui contient $\underline{O}.1_E$. Conditions équivalentes :
\begin{itemize}
  \item a) $A/\underline{O}1_E$ est monogène, i.e. $\exists$ une section $u$ de $A$ t.q. $A = \underline{O}.1_E + \underline{O}u$ ;
  \item a') $\exists$ loc. une section \ill{} ;
  \item b) $A$ est \uncertain{engendré} loc. par deux sections.
\end{itemize}
\marginal{résulte de \uncertain{N.C.} b) l'équation \ill{} $u^2 - \mathrm{Tr}(u)u + N(u) = 0$}
Un tel $A$ est stable par la loi d'algèbre de $\underline{\mathrm{End}}(E)$, et est un faisceau d'anneaux commutatif, et de plus \struck{entier} \add{loc.} \uncertain{fini} sur $\underline{O}$. Donc les $A$ en question correspondent : des \add{structures de} similitude sur $E$) correspondant \add{(commutatives)} aux sous-algèbres) de $\underline{\mathrm{End}}(E)$. Les str. de similitude str. fidèles correspondent loc. aux $A_u$ \struck{loc. définis par $u$} tels que $u$ ne soit

\page{15}
un scalaire sur aucune fibre résiduelle.

$A$ est appelé l'anneau des \emph{homothéties} (\add{sous-entendu :} \emph{directes}) \struck{pour} de la structure de similitude $\Sigma$, on le note $A_\Sigma$.

\emph{Lemme 2} Soit $\Sigma$ une structure de similitude sur $E$, et $A_\Sigma \subset \underline{\mathrm{End}}(E)$ l'Anneau des homothéties. Conditions équivalentes :
\begin{itemize}
  \item a) \struck{$\Sigma$ est} strictement fidèle \add{$A_\Sigma$ est loc. libre de rang 2 et}
  \item b) $A(s) \to \underline{\mathrm{End}}(E(s))$ injectif pour tout $s \in X$, i.e. $A$ sous-fibré vectoriel (de rang 2) de $\underline{\mathrm{End}}(E)$.
  \item b') (si $A = \underline{O}.1_E + \underline{O}.u$) $\forall s \in X$, $u(s) \in \underline{\mathrm{End}}(E(s))$ \add{n'est pas un scalaire}.
  \item c) $E$ est un $A$-Module \struck{loc. libre de rang} inversible, i.e. loc. libre de rang 1.
  \item c') $E$ est un $A$-module loc. monogène.
  \item d) $\mathcal{E} = \underline{\mathrm{End}}(E)$ est un $A$-Module ; c). Comme la structure régulière \ill{}) loc. libre de rang 2.
  \item d') $\mathcal{E}$ est un $A$-module loc. engendré par 2 éléments.
\end{itemize}
\note{au-dessus de d') : $\Delta \subset \underline{\mathrm{Quad}}(E)$ ; à la suite, « $x \in \Gamma E$, dans \ill{} » est biffé}

\emph{Cor. 1}, Sous ces conditions, soit $\Delta$ le sous-Module \struck{loc. lib.} inv. de $\underline{\mathrm{Quad}}(E)$ tel que $\Sigma = \Sigma_\Delta$. On a un iso. can.
\[
  \Delta \simeq (A/\underline{O}) \otimes L^\vee \qquad (\text{où } L = \textstyle\bigwedge^2 E = \det E)
\]
d'où
\[
  \check\Delta \simeq (A/\underline{O})^\vee \otimes L ,
\]

\page{16}
et on a \uncertain{la} forme quadratique de définition de la structure \struck{quadratique} \add{de similitude}
\[
  Q_\Sigma : E \longrightarrow \Delta^\vee \overset{\mathrm{df}}{=} \Delta' \quad (\simeq A/\underline{O} \otimes L^\vee)
\]
\struck{Soit pois, soit} Plus généralement, soit $\Delta'$ un Module inversible,
\[
  Q : E \longrightarrow \Delta' \qquad \bigl(\text{i.e. } \Delta'^\vee \to \underline{\mathrm{Quad}}(E) \text{ un homom. de Modules}\bigr)
\]
une forme quadratique, d'où une structure de similitude $\Sigma = \Sigma_Q$ sur $E$, correspondant à un sous-anneau $A = A_\Sigma$ de $\underline{\mathrm{End}}(E)$. Soit $x$ une section de $E$. Conditions équivalentes
\begin{itemize}
  \item a) $\{x\}$ est un générateur de $E$ comme $A$-Module, i.e. $A \to E$ \add{$a \mapsto a.x$}, est épi ;
  \item a') \struck{$Q(x)$ est} $\{x\}$ est une base de $E$ comme $A$-Module, i.e. $A \to E$ est un iso ;
  \item b) $\{Q(x)\}$ est une base de $\Delta'$, i.e. $Q(x) \in \Gamma \Delta'$ est inversible ;
  \item b') (si $\Delta' = \underline{O}$) $\forall s \in X$, $Q(x)(s) \neq 0$ ;
  \item \struck{c) Soit $Q$} \struck{Lem~2}
  \item d) (si $u$ est un gén. loc. de l'Alg.~$A$) $\forall s \in X$, $x(s) \in E(s)$ n'est pas \emph{vecteur propre} de $u(s) \in \underline{\mathrm{End}}\,E(s)$.
\end{itemize}

\page{17}
2.3. Considérons à nouveau l'iso. can.
\[
  \underline{\mathrm{End}}(E)/\underline{O}.1_E \simeq \underline{\mathrm{Quad}}(E) \otimes \underline{L} \qquad (\simeq \underline{\mathrm{Quad}}(E, L))
\]
\struck{Soit} $u$ une section de $\underline{\mathrm{End}}(E)$ i.e. un endom. de $E$, \struck{considérons} son image est donc une forme quadratique $Q_u$ à valeurs dans $L$. Je dis que
\[
  \boxed{\,Q_u(x) = u(x) \wedge x\,}
\]
(\struck{\ill{}} il faut vérifier le signe !). \emph{C'est} On vérifie directement par des coordonnées que
\[
  Q_u = 0 \iff u \in \Gamma\, \underline{O}.\mathrm{id}_E
\]

[NB. Si
\[
  u = \begin{pmatrix} a & b \\ c & d \end{pmatrix}
\]
alors
\[
  Q_u(\lambda, \mu) = c\lambda^2 + (d - a)\lambda\mu \ \ldots\ b\mu^2
\]
\note{le signe devant $b\mu^2$ est couvert d'une tache}
qui est identiquement nul en $\lambda, \mu$ ssi $b = c = 0$, $d = a$, i.e. $u = \begin{pmatrix} a & 0 \\ 0 & a \end{pmatrix} = a\,\mathrm{id}_E$.
Si
\[
  Q = \alpha\lambda^2 + \beta\lambda\mu + \gamma\mu^2 ,
\]
alors les $u$ tels que $Q_u = Q$ sont ceux de la forme
\[
  u_a = \begin{pmatrix} a & -\gamma \\ \alpha & a + \beta \end{pmatrix}, \qquad a \in \Gamma\,\underline{O} \text{ arbitraire},
\]
parmi lesquels
\[
  u_0 = \begin{pmatrix} 0 & -\gamma \\ \alpha & \beta \end{pmatrix} \quad \ldots \quad ]
\]

On a \struck{alors}

\page{18}
\[
  \boxed{\,\delta(u) \;\Bigl(\overset{\mathrm{df}}{=} (\mathrm{Tr}\,u)^2 - 4N(u)\Bigr) = -\underbrace{\delta(Q_u)}_{\text{discr. de la forme } Q_u}\,}
\]

— Attention au signe — dans le second cas !

On note que si $Q$ est une forme quadratique à valeurs dans le module inv. $\Delta'$ \struck{(\ill{} inv.)} \struck{inv./$\Delta$ $\Lambda$, alors}, alors
\[
  \delta(Q) \in \Gamma(\Delta'^{\otimes 2} \otimes L^{\otimes 2})
\]
est défini ainsi :
\[
  \varphi_Q : E \otimes E \longrightarrow \Delta' ,
\]
d'où
\[
  E \longrightarrow \underline{\mathrm{Hom}}(E, M) \simeq E^\vee \otimes \Delta' ,
\]
d'où \struck{\ill{}} en prenant l'ext. \uncertain{deux} \struck{dit} maximale, $\det = \bigwedge^2$ :
\[
  \underbrace{\det E}_{L} \longrightarrow \underbrace{\det(E^\vee)}_{L^\vee} \otimes \Delta'^{\otimes 2}
\]
d'où
\[
  \delta(Q) \in \Gamma\, \underbrace{L^{\vee \otimes 2} \Delta'^{\otimes 2}}_{(\Delta' \otimes L^\vee)^{\otimes 2}}
\]
\note{la formule encadrée plus haut écrivait $L^{\otimes 2}$ là où ce calcul donne $L^{\vee \otimes 2}$}

\struck{Mais d'autre part} Quand $\Delta' \simeq L$, c'est donc une section de $\underline{O}$, ce qui donne un sens à la formule sur les discriminants…

2.4. Soit
\[
  Q : E \longrightarrow \Delta' \qquad Q \in \Gamma\, \underline{\mathrm{Quad}}(E, \Delta')
\]

\page{19}
une forme quadratique à valeurs dans le module inversible $\Delta'$ \add{$= \Delta^\vee$}. Il revient au même de se donner
\[
  \Delta \longrightarrow \underline{\mathrm{Quad}}(E) \qquad (\simeq \underline{\mathrm{End}}(E)/\underline{O} \otimes L^\vee)
\]
(vrai pour $E$ loc. libre de r. fini), i.e. de se donner
\[
  u_Q : \Delta \otimes L \longrightarrow \underline{\mathrm{End}}(E)/\underline{O}.1
\]
Nous allons en profiter pour définir une alg. quadratique $A_Q$ dans $X$, et un homom. d'alg.
\[
  A_Q \longrightarrow \underline{\mathrm{End}}(E)
\]
dont l'image soit l'anneau $A_\Sigma$ des similitudes \add{$\Sigma = \Sigma_Q$} défini par la structure de similitude $\Sigma$ \add{i.e. l'image inverse de $\mathrm{Im}\,u_Q$ dans $\underline{\mathrm{End}}(E)$}, (pour \uncertain{ceci}), i.e. associée à la structure de $Q$. Considérons
\[
  0 \to \underline{O}_X \longrightarrow \underline{\mathrm{End}}(E) \longrightarrow \underline{\mathrm{End}}(E)/\underline{O}.1_E \to 0
\]
et \add{on désigne par $A_Q$} l'image inverse de cette extension par $u_Q$, d'où une extension

\begin{tikzcd}
  0 \arrow[r] & \underline{O} \arrow[r] \arrow[d, "\mathrm{id}"] & A_Q \arrow[r] \arrow[d, "u^!_Q"] & \Lambda \arrow[r] \arrow[d, "u_Q"] & 0 \\
  0 \arrow[r] & \underline{O} \arrow[r] & \underline{\mathrm{End}}(E) \arrow[r] & \underline{\mathrm{End}}(E)/\underline{O} \arrow[r] & 0
\end{tikzcd}

\note{au-dessus de $\Lambda$ : $\Lambda \overset{\mathrm{df}}{=} \Delta \otimes L$}

\page{20}
Considérons sur $A_Q$ la forme quadratique $N_Q$ définie par
\[
  N_Q(x) = N(u^!_Q(x))
\]
on aura donc, en désignant par $1_Q$ la section $1$ de $\underline{O}_X$, regardée comme section de $A_Q$,
\[
  N_Q(1_Q) = 1 ,
\]
Ainsi,
\[
  (A_Q, 1_Q, N_Q)
\]
est un module loc. libre de rang 2, muni d'une section $1_Q$ et d'une forme quadratique valant $1$ sur cette section. On a donc une structure d'alg. quadratique unique sur $A_Q$, pour laquelle $N_Q$ soit la norme. De plus, on vérifie immédiatement que $u^!_Q$ est un homom. d'algèbres, et que, en plus de la relation sur les normes, les relations
\[
  \mathrm{Tr}(u^!_Q(x)) = \mathrm{Tr}_{A_Q/\underline{O}}(x)
\]
\[
  u^!_Q(\bar x) = \overline{u^!_Q(x)}
\]
où on désigne par
\[
  v \longmapsto \bar v = (\mathrm{Tr}\,v).1 - v
\]
l'anti-involution canonique de $\underline{\mathrm{End}}(E)$, égale à $\mathrm{id}_{\underline{O}}$ sur $\underline{O}.1_E$ et induisant $-\mathrm{id}$ sur le quotient.
\note{l'argument se poursuit au-delà de la p.~20, hors de ce lot}

\end{document}
